\section{Discussion \& Conclusion}
\label{sec:conclusion}

\paragraph{Discussion.} Through a comprehensive review of existing research on the honesty of LLMs, we broadly categorize the honesty of LLMs into two key capabilities: self-knowledge and self-expression. Self-knowledge emphasizes the discriminative ability of LLMs to recognize and acknowledge the boundaries of their own knowledge, thereby mitigating undesirable behaviors such as hallucinations and overconfidence. In contrast, self-expression highlights the generative ability of LLMs to faithfully convey their inherent knowledge, whether derived from internal parameters or external context. Despite significant progress in this area, several key challenges remain that warrant further attention:
\begin{itemize}
\item Research on self-knowledge has primarily focused on short-form question-answering scenarios. However, it remains unclear how to enable an LLM to express its self-knowledge in real-world, long-form settings, particularly those that are complex and require a thoughtful reasoning process to identify its limitations.
\item Both self-knowledge and self-expression necessitate the identification of the model's inherent knowledge. Current approaches largely rely on intuitive methods, which may be suboptimal, such as determining knowledge based on its inclusion in the pre-training corpus or the model's ability to provide correct answers. More explainable identification methods, such as those based on hidden activations, are needed to ensure that the model's inherent knowledge is neither underestimated nor overestimated.
\item Current research on the honesty of LLMs mainly addresses textual inputs, while multimodal scenarios, including inputs with figures or videos, remain underexplored. It is unclear whether existing methods and findings can be effectively generalized to multimodal settings.
\end{itemize}

\paragraph{Conclusion.} Honesty is a crucial factor in the development of LLMs, yet current models still exhibit significant dishonest behaviors. To address these issues, this paper offers a thorough overview of research on the honesty of LLMs, including its clarification, evaluation approaches, and improvement strategies. Furthermore, we propose several potential directions for future research. We hope this survey serves as a valuable resource for researchers studying LLM honesty and encourages further exploration in this field.
